% Part: incompleteness
% Chapter: arithmetization-syntax
% Section: introduction

\documentclass[../../../include/open-logic-section]{subfiles}

\begin{document}

\olfileid{inc}{art}{int}
\olsection{مقدمة}

نحتاج، للربط بين قابلية الحساب والمنطق، إلى طريقة للحديث عن كائنات
المنطق (الرموز والحدود وال!!{formula}s وال!!{derivation}s)، وعن
العمليات عليها وخواصها وعلاقاتها، على نحو ملائم للمعالجة الحسابية.
ويمكننا فعل ذلك مباشرة بالنظر في الدوال والعلاقات القابلة للحساب على
الرموز ومتتاليات الرموز وسائر الكائنات المبنية منها. ولما كانت كائنات
التركيب المنطقي كلها منتهية ومبنية من مجموعات !!a{enumerable} من
الرموز، كان ذلك ممكنًا في بعض نماذج الحساب. لكن نماذج أخرى للحساب،
مثل الدوال العودية، مقصورة على الأعداد وعلاقاتها ودوالها. وفوق ذلك
نريد في نهاية المطاف أن نستطيع معالجة التركيب داخل نظريات معينة،
وبخاصة النظريات المصوغة في لغة الحساب. وفي هذه الحالات لا بد من
\emph{حوسبة} التركيب؛ أي تمثيل الكائنات التركيبية والعمليات عليها
وعلاقاتها، على الترتيب، بأعداد ودوال حسابية وعلاقات حسابية. والفكرة،
التي ترجع إلى لايبنتس، هي إسناد أعداد إلى الكائنات التركيبية.

من السهل نسبيًا إسناد أعداد إلى الرموز بوصفها «رموزها العددية». وتشكل
بعض الرموز تحديًا صغيرًا؛ فثمة مثلًا !!{variable}s غير متناهية، بل
!!{function}s غير متناهية من كل رتبة~$n$. لكن يمكن بالطبع إسناد
الأعداد إلى الرموز منهجيًا بحيث يسند إلى $\Obj v_2$ و$\Obj v_3$،
مثلًا، رمزان عدديان مختلفان. أما متتاليات الرموز (كالحدود
وال!!{formula}s) فتشكل تحديًا أكبر. لكن إذا استطعنا معالجة متتاليات
الأعداد حسابيًا محضًا (بترميز المتتاليات بقوى الأعداد الأولية مثلًا)،
أمكننا توسيع ترميز الرموز المفردة إلى ترميز متتاليات الرموز، ثم إلى
متتاليات ال!!{formula}s أو ترتيبات أخرى منها، كال!!{derivation}s.
ويسمى هذا الترميز الموسع «ترقيم غودل». ويسند إلى كل حد و!!{formula}
و!!{derivation} عدد غودل.

بترميز متتاليات الرموز بمتتاليات رموزها العددية، واختيار نظام لترميز
المتتاليات يمكن معالجته بدوال قابلة للحساب، نستطيع أيضًا معالجة أعداد
غودل بدوال قابلة للحساب. وفي التطبيق العملي ستكون جميع الدوال المعنية
عودية بدائية. فمثلًا، حساب طول متتالية وحساب عنصرها رقم $i$ من رمزها
العددي كلاهما عودي بدائي. وإذا كان العدد الذي يرمز المتتالية، مثلًا،
عدد غودل ل!!a{formula}~$!A$، رأينا فورًا أن طول !!a{formula} والرمز
العددي للرمز رقم $i$ في !!a{formula} يمكن حسابهما أيضًا من عدد غودل لـ~$!A$.
ومن الأصعب قليلًا إثبات أن خاصية كون عدد ما عدد غودل لحد حسن التكوين
أو ل!!{derivation} صحيح عودية بدائية. ومع ذلك فهذا ممكن، لأن
المتتاليات المعنية (الحدود وال!!{formula}s وال!!{derivation}s) معرفة
استقرائيًا.

انظر، مثالًا، إلى عملية الاستبدال. فإذا كانت $!A$ صيغة، وكان $x$
متغيرًا و$t$ حدًا، فإن $\Subst{!A}{t}{x}$ نتيجة استبدال~$t$ بكل ورود
حر لـ~$x$ في~$!A$. وافترض الآن أننا أسندنا أعداد غودل إلى $!A$ و$x$
و$t$، ولتكن $k$ و$l$ و$m$، على الترتيب. ويسند المخطط نفسه عدد غودل
إلى $\Subst{!A}{t}{x}$، وليكن~$n$. والتطبيق الذي يرسل $k$ و$l$ و$m$
إلى $n$ هو النظير الحسابي لعملية الاستبدال. فحين ترسل عملية الاستبدال
$!A$ و$x$ و$t$ إلى $\Subst{!A}{t}{x}$، ترسل دالة الاستبدال المحوسبة
أعداد غودل $k$ و$l$ و$m$ إلى عدد غودل~$n$. وسنرى أن هذه الدالة
عودية بدائية.

ليست حوسبة التركيب ذات أهمية مجردة فحسب، وإن كان اكتشاف أن لغات مثل
لغة الحساب، التي لا تأتي بآليات «للحديث عن» اللغات، تستطيع مع ذلك
صوغ خواص معقدة للتعبيرات صوغًا صوريًا، اكتشافًا مهمًا غير تافه. ولا
يبقى بعد ذلك إلا خطوة صغيرة لنسأل عما تستطيع نظرية بهذه اللغة، مثل
حساب بيانو، أن \emph{تثبته} عن لغتها نفسها (ومن ذلك، مثلًا، ما إذا
كانت ال!!{sentence}s قابلة للبرهنة أو صادقة). وهذا يقودنا إلى
مبرهنات الحدود الشهيرة لغودل (في عدم قابلية البرهنة) وتارسكي (في عدم
قابلية تعريف الصدق). لكن حيلة حوسبة التركيب مهمة أيضًا لإثبات بعض
النتائج المهمة في نظرية قابلية الحساب، مثل النتائج المتعلقة بالقدرة
الحسابية للنظريات أو بالعلاقة بين نماذج قابلية الحساب المختلفة.
وتصلح حوسبة التركيب نموذجًا لحوسبة كائنات وخواص أخرى. فمن الممكن،
مثلًا، حوسبة التهيئات والحسابات (لآلات تورنغ مثلًا) على نحو مماثل.
ويتيح ذلك محاكاة الحسابات في نموذج (كآلات تورنغ) داخل نموذج آخر
(كالدوال العودية).

\end{document}
